\ssscp{Putative PCEMP *o}{13-01-05-01}
In \citet[247--249]{Blust1993} the evidence for PCEMP *o is based on
two purported lexical developments: PMP *tudan > PCEMP *todan
‘sit’ and PMP *uliq > PCEMP *oliq ‘return’.

With regards to purported PMP *tudan > PCEMP *todan ‘sit’,
the geographical distribution of forms in \citet{Blust1993} supporting
*todan is limited to a handful of languages clustered near the
EMP boundary: five languages from our \tsc{Seram-Tanimbar-Bomberai}
subgroup (\ili{Geser}, \ili{Kowiai}, \ili{Yamdena}, \ili{Sekar}, \ili{Irarutu}),
and two SHWNG languages (\ili{Buli},
Minyaifuin) from outside our target region.
Languages such as \ili{Yamdena},
\ili{Geser}, and \ili{Kowiai} have many words with \it{o} or \it{u} regularly
resulting from PMP *u.
\ili{Buli} and many other SHWNG languages have a 7-vowel system,
so the presence of \it{o} is expected.
No supporting evidence is given from languages belonging to \tsc{\ili{Bima}-Lembata},
\tsc{Timor-Babar}, \tsc{\ili{Central} Timor}, \tsc{Aru},
\tsc{Sula-Buru}, or \tsc{\ili{Ambon}-Seram} subgroups.

After responding to arguments
and objections from \citet{Nothofer1992}
and \citet{DonohueGrimes2008},
\citet[57]{Blust2009} says, “I will abandon my claim that it
[*todan ‘sit’] provides support for CEMP.” \citet{BlustTrussel2020}
give an updated form for PMP *tu(n)daŋ ‘sit’ based on only six
languages in total, only two of which have /o/.
The STB language \ili{Geser} \it{ma-toran} ‘sit’
and SHWNG \ili{Buli} \it{totolaŋ} have a note added “with irregular
lowering of *u to \it{o} in the PCEMP languages \ili{Geser}
and \ili{Buli}.” It may not be so irregular in \ili{Geser}.
In a quick scan of \ili{Geser} data,
we found seven other words in which PMP *u lowered to \it{o},
indicating a split *u > \it{u,
o} commonly found in other languages across the region.\footnote{
		For
		\ili{Geser}: *babuy > \it{boi} ‘pig’,
		*timuR > \it{timor-a} ‘\ili{east}’,
		*baqəRu > \it{wowo} ‘new’,
		*qatəluR > \it{tolor} ‘egg’,
		*təlu > \it{tolo} ‘three’,
		*manuk ‘bird’ > \it{manok takok} ‘chicken’,
		*quRis (\tsc{PPh}, \citealp{ZorcCharles1971}) ‘clean,
		scrape’ > \it{ori} ‘clean,
		scrape for use’ (cf.
		also \ili{Buru} \it{ohi} ‘clean (rattan),
		scrape away from the body’; \tsc{POc} *ori(t) ‘scratch, scrape, peel’).
		In the mainland \ili{Waru} dialect of \ili{Geser},
		some of these are conditioned by harmony of vowel height where
		*ə > \it{o} in the penultimate syllable.}

With regards to putative PMP *uliq > PCEMP *oliq ‘return’,
\citet{Blust2009} and \citet{BlustTrussel2020} acknowledge
that in the CEMP region,
*uliq and *oliq coexist side-by-side in the same subgroups,
so it is not a replacement innovation.\footnote{
In \citet[410f]{Ross2016},
putative \tsc{POc} *oli(q) is reflected in only one subgroup,
\tsc{South\ili{east} Solomonic}. Malcolm Ross (pers. comm.
September 2022) says he would now see this as a “local idiosyncratic
lexical innovation”, and also observes that \citet{Kamholz2014} does not report any
SHWNG reflexes for *oliq.}

Forms with \it{u}, such as \tsc{\ili{Ambon}-Seram} language South \ili{Nuaulu}
\it{n-uni} ‘3s-return’ (*l > \it{n} is regular),
\tsc{Seram-Tanimbar-Bomberai} languages \ili{Masiwang} \it{uli} ‘return,
go back’, \ili{Geser} \it{muli} ‘return (home)’,
SHWNG language \ili{Gane}/\ili{Giman} \it{mul} ‘return home’,
are found in various subgroups throughout the CEMP region,
extending out to Oceania as far as \ili{Rarotongan}.\footnote{
It is
unclear if the \ili{Geser}
and \ili{Gane}/\ili{Giman} forms with initial {\#}m (taken from word lists),
reflect a \tsc{2sg} prefix \it{m-},
or the PMP actor voice infix *\<um\>,
which then became a prefix before a vowel-initial root.}

Outside the target region \citet[879]{Mills1975} reconstructs
Proto-South Sulawesi *\it{-ole} ‘again,
return’ (Proto-South Sulawesi *-e{\#} is the regular correspondence for PMP *-iq{\#}).
\citet[247]{Blust1993} dismisses this reconstruction saying “Although
Mills actually reconstructs PSS [Proto-South Sulawesi] *p/ole,
*pa/ole, and *m/ole, *ma/ole,
the supporting evidence is consistent with *pa-ule,
*ma-ule, with the first-syllable mid-vowel in many reflexes a
product of contraction across a morpheme boundary.” The surface
forms in the South Sulawesi languages have \it{o},
and Proto-South Sulawesi is reconstructed with mid vowels *o
and *e, as noted earlier.
So Blust chooses to interpret the South Sulawesi data differently
than Mills and Nothofer.
Yet \citet[247]{Blust1993} acknowledges “This leaves only \ili{Muna}
\it{d-oli} 'turn' as a possible counter-example to the proposed
PCEMP innovation.” 
Only six non-Oceanic “CMP” area languages are given in \citet{Blust1993}
in support of PCEMP *oliq ‘return’.
We add more languages in this discussion.
In \tsc{Sula-Buru}, PMP *u underwent an unconditioned split >
\it{u, o} in all \tsc{Sula-Buru} languages.
So \ili{Buru} \it{oli} ‘return,
come back; marry, revert to’
and \it{epyolik} /ep-oli-k/ ‘return something,
bring something back’ \citep{GrimesB1990,GrimesGrimes2020} can be derived directly
from PMP *uliq ‘return’ through normal sound changes.
The proposed PCEMP form is not needed to explain the \ili{Buru} form.

The \ili{Selaru} (\tsc{Timor-Babar}) of southwest Tanimbar
has \it{olik} ‘return, untangle’.
There is a split in PMP *u > \it{u,
o, w}, so the \ili{Selaru} form can also derive directly from the PMP form.
A perusal of PMP words with *u shows that the processes in \ili{Selaru}
are independent from the processes in \ili{Buru} (in addition to being
in different subgroups) and that sometimes they align.
For example, PMP *upu,
*əmpu ‘grandparent, grandchild, master, lord’ > \ili{Buru} \it{opo},
\ili{Selaru} \it{ebu}; PMP *qulu ‘head,
leader’ > \ili{Buru} \it{olo} ‘head of social group’,
\ili{Selaru} \it{uskwe} ‘head, chief’ (/usu-ke/,
regular split of PMP *l > \it{l,
s}); PMP *qahəlu ‘pestle’ > \ili{Buru} \it{alu},
\ili{Selaru} \it{asw}; PMP *suqan ‘digging stick’ > \ili{Buru} \it{suan},
\ili{Selaru} \it{su}; PCEMP *upi ‘blow’ > \ili{Buru} \it{opi},
\ili{Selaru} \it{ofw {\tl} ohw}.
So PMP *uliq > \ili{Buru} \it{oli}
and \ili{Selaru} \it{olik} are from parallel development,
not through shared inheritance.

The \tsc{Seram-Tanimbar-Bomberai} languages \ili{Yamdena}
and \ili{Fordata} terms “diverge semantically” from PMP *uliq ‘return’.
Compare \ili{Yamdena} \it{n-ol, n-oli} ‘\tsc{3sg}-to various places’;\footnote{
		\citet[247]{Blust1993}
		glosses \ili{Yamdena} \it{n-ol} as ‘move to and fro’.
		The gloss from \citet{MettlerMettler1996} as ‘ke mana-mana (to
		various places)’ is more consistent with the gloss of the cognate
		term in \ili{Fordata} from an independent source.}
\ili{Fordata} \it{n-oli} ‘\tsc{3sg}-walk around’.
(\ili{Yamdena} also has a motion verb \it{bali} + pronoun ‘return’
< PMP *balik ‘reverse,
turn around’; \ili{Fordata} has \it{n-ewal} ‘return’.) \ili{Yamdena} splits
*u > \it{u, o, i/y}{\#}, so again parallel development is a reasonable explanation.\footnote{
		\tsc{Southwest
		Maluku} languages \ili{Luang} \it{wali} ‘return home’; \ili{Kisar} \it{wali}
		‘return’ are more likely to be from PMP *balik ‘reverse,
		turn around’, than from either *uliq or *oliq.
		Compare other evidence for *b > \it{w} in PMP *baliw ‘counterpart’
		> \ili{Luang} \it{ha-wali} ‘boyfriend,
		girlfriend’; PMP *bulu ‘body hair,
		feather, fur’ > \ili{Luang} \it{wulu} ‘body hair,
		feather, fur’; PMP *balu ‘widow’ > \ili{Kisar} \it{wal{\tl}walum} ‘widow’;
		PMP *buku ‘joint’ > \ili{Kisar} \it{wuʔu-n} ‘knot,
		joint’.} 
The unpatterned distribution of \it{o}-forms interspersed alongside
cognates with \it{u} is found repeatedly in the region in Austronesian
vocabulary as shown in examples \qf{13-05}--\qf{13-08}.
To put it another way,
a similar pattern is found widely in many etyma in many languages
in the region, not just in *oliq.
Forms with \it{o} are shaded.
The data are so inconsistent as to which languages have \it{u}
or \it{o} in which etyma,
as to not be amenable to combining in tables with the same languages.

{\st{6pt}\begin{xltabular}[l]{\linewidth}
{>{\hspace{-\tabcolsep}}l<{\hspace{-\tabcolsep}}
>{\hspace{-\tabcolsep}\enspace}lll
>{\raggedright\arraybackslash}X}
\lsptoprule
\tbx{13-05}	&	PMP	&	*qulu	&	head	&	(cognates with \it{o} are also found in Sulawesi)	\\	\midrule
\endfirsthead
\lsptoprule
\qf{13-05}	&	PMP	&	*qulu	&	head	&		\\	\midrule
\endhead
	&	\ili{Manggarai}	&	ulu	&	head	&	\tsc{BL}	\\
	&	\ili{Ngadha}	&	ulu	&	head	&	\tsc{BL}	\\
\rcl	&	\ili{Lio}	&	holo	&	head	&	\tsc{BL}	\\
	&	\ili{Ende}	&	uru	&	head	&	\tsc{BL}	\\
	&	S. \ili{Mambae}	&	ulu hatu	&	head	&	\tsc{CT}	\\
	&	\ili{Tetun}	&	ulu-n	&	head	&	\tsc{TB}	\\
	&	\ili{Makatian}	&	ʔuʤun-e	&	head	&	\tsc{TB}	\\
	&	\ili{Selaru}	&	usu khatu	&	head	&	\tsc{TB}, \tsc{SW Tanimbar}	\\
	&	\ili{Dobel}	&	kwulu-	&	head	&	\tsc{Aru}	\\
	&	\ili{Ujir}	&	ulu-	&	head	&	\tsc{Aru}	\\
	&	\ili{Kei} Kecil	&	uun	&	head, brain, source	&	\tsc{STB}	\\
	&	\ili{Watubela}	&	kulu	&	head	&	\tsc{STB}	\\
	&	\ili{Banda}	&	ulu-n	&	head	&	\ili{Banda}	\\
	&	\ili{Liana-Seti}	&	una	&	head	&	\tsc{AS}	\\
	&	\ili{Alune}	&	ulu-i	&	head	&	\tsc{AS}	\\
\rcl	&	\ili{Larike}	&	udo hatu	&	head	&	\tsc{AS}	\\
	&	\ili{Asilulu}	&	ulu-	&	head	&	\tsc{AS}	\\
	&	Boano	&	unu	&	head	&	\tsc{AS}	\\
\rcl	&	\ili{Kayeli}	&	olo-ni	&	head	&	\tsc{AS}	\\
\rcl	&	\ili{Ambelau}	&	olo-ni	&	head	&	\tsc{SB}	\\
\rcl	&	\ili{Buru}	&	olo, olo-n	&	head	&	\tsc{SB}	\\
	&	\ili{Lisela}	&	ulu-n	&	head	&	\tsc{SB}	\\
\lspbottomrule
\end{xltabular}\addtocounter{table}{-1}}

{\st{6pt}\begin{xltabular}[l]{\linewidth}
{>{\hspace{-\tabcolsep}}l<{\hspace{-\tabcolsep}}
>{\hspace{-\tabcolsep}\enspace}lll
>{\raggedright\arraybackslash}X}
\lsptoprule
\tbx{13-06}	&	PMP	&	*bulu	&	hair, fur, feather	&	(cognates with \it{o} also in Borneo and \ili{Malagasy})	\\	\midrule
\endfirsthead
\lsptoprule
\qf{13-06}	&	PMP	&	*bulu	&	hair, fur, feather	&		\\	\midrule
\endhead
	&	\ili{Manggarai}	&	ʋulu	&	feather	&	\tsc{BL}	\\
	&	So{\Q}a	&	fuu	&	feather	&	\tsc{BL}	\\
	&	\ili{Kambera}	&	wulu	&	feather	&	\tsc{BL}	\\
	&	\ili{Tetun}	&	fulu-n	&	body hair	&	\tsc{TB}	\\
	&	\ili{Galolen}	&	hulu-r	&	body hair	&	\tsc{TB}	\\
\rcl	&	\ili{Batuley}	&	ɸoloi	&	feather, body hair	&	\tsc{Aru}	\\
	&	\ili{Fordata}	&	vulu-n	&	body hair	&	\tsc{STB}	\\
	&	\ili{Kei} Kecil	&	vuu-n	&	body hair	&	\tsc{STB}	\\
	&	\ili{Kowiai}	&	ɸur	&	body hair	&	\tsc{STB}	\\
	&	\ili{Geser}	&	ulu	&	body hair	&	\tsc{STB} (*bu > \it{u})	\\
	&	\ili{Masiwang}	&	hul-huli-n	&	body hair	&	\tsc{STB}	\\
	&	\ili{Tamilou}	&	huyu-n	&	body hair	&	\tsc{AS}	\\
	&	\ili{Alune}	&	bulu-i	&	hair, fur, feather	&	\tsc{AS}	\\
	&	\ili{Larike}	&	hudu-na	&	body hair	&	\tsc{AS}	\\
	&	\ili{Asilulu}	&	hulu-	&	body hair	&	\tsc{AS}	\\
	&	Boano	&	hunu	&	body hair	&	\tsc{AS}	\\
\rcl	&	\ili{Kayeli}	&	bolo-n	&	hair, fur, feather	&	\tsc{AS}	\\
\rcl	&	\ili{Ambelau}	&	bolu-e	&	body hair	&	\tsc{SB}	\\
\rcl	&	\ili{Buru}	&	folo-n	&	hair, fur, feather	&	\tsc{SB}	\\
	&	\ili{Lisela}	&	ɸulu-n	&	hair, fur, feather	&	\tsc{SB}	\\
\rcl	&	\ili{Sanana}	&	foa	&	body hair	&	\tsc{SB}	\\
\rcl	&	\ili{Mangole}	&	foo	&	body hair	&	\tsc{SB}	\\
\lspbottomrule
\end{xltabular}\addtocounter{table}{-1}}

{\st{6pt}\begin{xltabular}[l]{\linewidth}
{>{\hspace{-\tabcolsep}}l<{\hspace{-\tabcolsep}}
>{\hspace{-\tabcolsep}\enspace}lll
>{\raggedright\arraybackslash}X}
\lsptoprule
\tbx{13-07}	&	PMP	&	*qatəluR	&	egg	&	(cognates with final \it{o} are also found in western Indonesia and Sulawesi)	\\	\midrule
\endfirsthead
\lsptoprule
\qf{13-07}	&	PMP	&	*qatəluR	&	egg	&	\\	\midrule
\endhead
	&	\ili{Bima}	&	dolu	&	egg	&	\tsc{BL}	\\
\rcl	&	\ili{Manggarai}	&	təlo	&	egg	&	\tsc{BL}	\\
\rcl	&	\ili{Ngadha}	&	təlo	&	egg	&	\tsc{BL}	\\
\rcl	&	\ili{Lio}	&	təlo	&	egg	&	\tsc{BL}	\\
\rcl	&	Palu{\Q}e	&	təlo	&	egg	&	\tsc{BL}	\\
\rcl	&	\ili{Sika}	&	təlo	&	egg	&	\tsc{BL}	\\
	&	\ili{Lamalera}	&	təlu	&	egg	&	\tsc{BL}	\\
\rcl	&	\ili{Kedang}	&	tolor	&	egg	&	\tsc{BL}	\\
	&	\ili{Anakalang}	&	tilu	&	egg	&	\tsc{BL}	\\
	&	\ili{Wanukaka}	&	tilu	&	egg	&	\tsc{BL}	\\
	&	\ili{Laboya}	&	talu	&	egg	&	\tsc{BL}	\\
	&	\ili{Havu}	&	dəlu	&	egg	&	\tsc{BL}	\\
\rcl	&	\ili{Kotos} \ili{Amarasi}	&	teno-ʔ	&	egg	&	\tsc{TB}, \tsc{Meto}	\\
\rcl	&	\ili{Dela}	&	manu telo-ʔ	&	egg	&	\tsc{TB}, \tsc{W. Rote-Meto}	\\
\rcl	&	\ili{Termanu}	&	tolo	&	egg	&	\tsc{TB}, \tsc{Nuclear Rote}	\\
\rcl	&	\ili{Rikou}	&	man toloʔ	&	egg	&	\tsc{TB}, \tsc{Nuclear Rote}	\\
	&	\ili{Tetun}	&	tolun	&	egg	&	\tsc{TB}	\\
	&	\ili{Galolen}	&	manu telun	&	egg	&	\tsc{TB}	\\
	&	\ili{Raklungu}	&	manuʔ teluʔ	&	egg	&	\tsc{TB}	\\
	&	\ili{Tugun}	&	telu-n	&	egg	&	\tsc{TB}	\\
	&	\ili{Kisar}	&	keru-n	&	egg	&	\tsc{TB} (regular *t > \it{k})	\\
	&	\ili{Selaru}	&	tesu	&	egg	&	\tsc{TB}	\\
	&	\ili{Makatian}	&	teʤu-n-e	&	egg	&	\tsc{TB}	\\
\rcl	&	\ili{Kemak}	&	man telo-n	&	egg	&	\tsc{CT}	\\
\rcl	&	S. \ili{Mambae}	&	maun telo-n	&	egg	&	\tsc{CT}	\\
	&	\ili{Kola}	&	tulih	&	egg	&	\tsc{Aru}	\\
	&	\ili{Ngaibor}	&	tulir	&	egg	&	\tsc{Aru}	\\
	&	\ili{Irarutu}	&	trú	&	egg	&	\tsc{STB}	\\
	&	\ili{Fordata}	&	mata telur	&	egg	&	\tsc{STB}	\\
	&	\ili{Kei} Kecil	&	tilur	&	egg	&	\tsc{STB}	\\
	&	\ili{Watubela}	&	katlu	&	egg	&	\tsc{STB}	\\
\rcl	&	\ili{Kowiai}	&	toro-n	&	egg	&	\tsc{STB}	\\
\rcl	&	\ili{Geser} (\ili{Waru})	&	tolor	&	egg	&	\tsc{STB}	\\
	&	\ili{Masiwang}	&	toli-n	&	egg	&	\tsc{STB}	\\
	&	\ili{Banda}	&	tuluru	&	egg	&	\ili{Banda}	\\
	&	\ili{Liana-Seti}	&	tolla	&	egg	&	\tsc{AS}	\\
\rcl	&	\ili{Manusela}	&	tolo-nia	&	egg	&	\tsc{AS}	\\
	&	\ili{Tehoru}	&	tei-n	&	egg	&	\tsc{AS}	\\
	&	\ili{Sepa}	&	toru-n	&	egg	&	\tsc{AS}	\\
	&	\ili{Yalahatan}	&	telul-i	&	egg	&	\tsc{AS}	\\
	&	S. \ili{Nuaulu}	&	toun-e	&	egg	&	\tsc{AS}	\\
	&	\ili{Alune}	&	manu telu-i	&	egg	&	\tsc{AS}	\\
	&	\ili{Haruku}	&	man teru	&	egg	&	\tsc{AS}	\\
	&	\ili{Asilulu}	&	man telu	&	egg	&	\tsc{AS}	\\
\rcl	&	\ili{Larike}	&	man tido	&	egg	&	\tsc{AS}	\\
\rcl	&	\ili{Alang}	&	man tilo-nu	&	egg	&	\tsc{AS}	\\
\rcl	&	\ili{Ambelau}	&	roho-i	&	egg	&	\tsc{SB} (reg. *t > \it{r}; *l > \it{l, h})	\\
	&	\ili{Buru}	&	teput telu-n	&	chicken egg	&	\tsc{SB}	\\
	&	\ili{Mangole}	&	manu telu	&	bird egg	&	\tsc{SB}	\\
\lspbottomrule
\end{xltabular}\addtocounter{table}{-1}}

\newpage
{\st{6pt}\begin{xltabular}[l]{\linewidth}
{>{\hspace{-\tabcolsep}}l<{\hspace{-\tabcolsep}}
>{\hspace{-\tabcolsep}\enspace}lll
>{\raggedright\arraybackslash}X}
\lsptoprule
\tbx{13-08}	&	PMP	&	*ma-qudip{\footnotemark}  	&	alive, live 	&	(\it{o} also in Taiwan, Philippines and western Indonesia)		\\	\midrule
\endfirsthead
\lsptoprule
	&	PMP	&	*ma-qudip  	&	alive, live 	&\\	\midrule
\endhead
\rcl	&	\ili{Bima}	&	mori	&	alive, live	&	\tsc{BL}	\\
	&	Ngad'a	&	muzi	&	alive	&	\tsc{BL}	\\
	&	\ili{Anakalang}	&	múruk	&	alive	&	\tsc{BL}	\\
\rcl	&	\ili{Wanukaka}	&	morik	&	alive	&	\tsc{BL}	\\
\rcl	&	Lamboya	&	morha	&	alive	&	\tsc{BL}	\\
	&	\ili{Dhao}	&	muri	&	live, grow, develop	&	\tsc{BL}	\\
\rcl	& \ili{Kotos}	\ili{Amarasi}	&	n-moni	&	live	&	\tsc{TB}	\\
\rcl	&	\ili{Tii}	&	horis	&	alive, live	&	\tsc{TB}	\\
\rcl	&	\ili{Termanu}	&	moli	&	live, grow, come up	&	\tsc{TB}	\\
\rcl	&	\ili{Tetun}	&	moris	&	live, alive	&	\tsc{TB}	\\
\rcl	&	\ili{Galolen}	&	mori	&	live, alive	&	\tsc{TB}	\\
\rcl	&	\ili{Raklungu}	&	lik-mori	&	alive, live	&	\tsc{TB}	\\
\rcl	&	\ili{Tugun}	&	mori	&	alive	&	\tsc{TB}	\\
\rcl	&	\ili{Kisar}	&	mori	&	alive	&	\tsc{TB}	\\
\rcl	&	\ili{Luang}	&	mori	&	alive, live, give birth	&	\tsc{TB}	\\
\rcl	&	\ili{Selaru}	&	morif	&	alive	&	\tsc{TB}	\\
\rcl	&	S. \ili{Mambae}	&	mori	&	alive, live	&	\tsc{CT}	\\
\rcl	&	\ili{Yamdena}	&	n-morip	&	live	&	STB	\\
\lspbottomrule
\end{xltabular}\addtocounter{table}{-1}}
\footnotetext{Forms with \it{o} are fairly common for this etymon,
		because of the widespread coalescence of *aqu/*a(C)u/*au
		> \it{o} in the region.
		So *ma-qudip > **maudip with subsequent **au > \it{o}.}

Like \ili{Buru}
and \ili{Selaru}, \tsc{\ili{Bima}-Lembata} languages \ili{Manggarai} \it{ole} ‘restore,
repair’ \ili{Havu} \it{pe-oli} ‘back-and-forth aimlessly,
to-and-fro, tack back-and-forth (sailing)’,\footnote{
		Our thanks to Maria Margereta (Nadope) Ga Kore on Sabu for clarifying the
		meaning of this \ili{Havu} form.} 
and \ili{Lamaholot} \it{oli-n} ‘return’ are more likely to have an
\it{o}-vowel through parallel development than through shared inheritance.
For example, regarding \ili{Manggarai},
\citet[89]{Blust2008} assigns \ili{Manggarai} \it{ole} ‘restore,
repair’ to PMP *uliq,
and says, “PMP *u usually became \it{u} [{\ldots}] occasionally,
however, it became \it{o}”, and provides six additional examples
from \ili{Manggarai}, including *uak > \it{oak} ‘nausea,
retching’, *ukuk > \it{okok {\tl} okuk} ‘sound of coughing’,
*pənuq > \it{peno} ‘full’
and *tuzuq > \it{toso} ‘point out’.

Methodologically, \citet[53]{Blust2009} says, “For the change
*uliq > *oliq to carry weight as subgrouping evidence,
all that matters is that convergence can be ruled out as a plausible
explanation for the shared innovation,
and that this innovation is exclusively shared by the languages
in the proposed subgroup.” Exactly.
Convergence cannot be ruled out since it is the more reasonable
explanation, because a similar pattern is found widely in other etyma in many
languages in the region,
and not just in this one etymon.
It is not an exclusively shared innovation,
and its distribution within the target region is not patterned
in ways useful for subgrouping.

Elsewhere \citep[42]{Blust2009} PCEMP *kandoRa ‘cuscus’ is given
as providing additional support for PCEMP *o.
Yet outside of EMP,
\citet{BlustTrussel2020} have only one (!) witness that supports
the full reconstruction, and that is right at the EMP boundary:
\ili{Watubela} \it{kadola} ‘marsupial’.
South Halmahera witnesses are ambiguous: \ili{Buli} \it{do} ‘kind of
small marsupial’, and \ili{Taba} \it{do} ‘cuscus’.
In \srf{14-03-04} we discuss several additional problems associated
with the anomalous \ili{Watubela} form,
its vowels and alternate sources for the SHWNG forms that better
reflect more languages in the region.

We note that all historical reconstructions listed with word
initial *o in \citet{BlustTrussel2020} are, without exception, \tsc{POc}.
None are attributed to PCEMP.

To recap, the evidence for PCEMP *o was originally based on developments
in two words, to which a third was added.
Blust has since abandoned his claim based on *todan ‘sit’.
*kandoRa ‘cuscus’ has only one witness outside of EMP right at
the boundary, and has multiple questions tied to that etymon,
its gloss and its distribution.
Given that a similar pattern happens in many words in the region -- some
with \it{u} and some with \it{o},
parallel development is a more reasonable explanation for reflexes
with \it{o} from *uliq ‘return’ than is attributing those forms
to an irregular sound change from a common parent language or linkage.